\section{Метод решения}
Код данной лабораторной работы был написан так, чтобы его можно было запускать как на операционных системах POSIX, так и на Windows.
Это было сделано благодаря реализации общего интерфейса для взаимодействия с операционной системой.
При запуске основной программы помимо родительского процесса запускаются два дочерних, все процессы связаны между собой через 3 pipe'a, с системным вводом и выводом взаимодейтсвует только родительский процесс, ввод и вывод дочерних процессов перенаправляются в pipe'ы

\section{Описание программы}
Общий интерфейс взаимодействия с операционной системой описан в заголовочном файле src/os\_interaction/basic/os.h.
При компиляции программы в ней выбирается та или иная реализация требуемых в работе функций, файлы с реализацией этих функций под разные системы находятся в папках src/os\_interaction/POSIX и src/os\_interaction/Windows
Основной код программы находится в файлах main.cpp, child1.cpp, child2.cpp

\textit{Осонвные структуры}
\begin{lstlisting}[language=C++]
struct FileHandleDummy { };
typedef FileHandleDummy* FileHandle;

struct ProcessHandleDummy { };
typedef ProcessHandleDummy* ProcessHandle;
\end{lstlisting}
Две структуры для описания HANDLE файлов и процессов для Windows.

\begin{lstlisting}[language=C++]
struct Process {
    int pid = -1;
    ProcessHandle handle = nullptr;
};

struct File {
    int id = -1;
    FileHandle handle = nullptr;
};
\end{lstlisting}
Базовые структуры для описания процессов в Windows и POSIX.

\begin{lstlisting}[language=C++]
struct ProcessStartInfo {

    std::string cmd;
    std::vector<std::string> args;

    File redirectStdin = {-1,nullptr};
    File redirectStdout = {-1,nullptr};

    ProcessStartInfo() = default;
};
\end{lstlisting}
Структура для описания параметров инициализации процесса (ключевые параметры - cmd, redirectStdin и redirectStdout).

\textit{Используемые системные вызовы}

Windows:
\begin{itemize}
  \item CreateProcess - создание процесса.
  \item TerminateProcess - принудительная остановка процесса.
  \item CreatePipe - создание pipe.
  \item GetStdHandle - получение стандартного потока ввода и вывода.
  \item Write - запись в поток.
  \item Read - чтение из потока.
\end{itemize}

POSIX:
\begin{itemize}
  \item posix\_spawn - создание процесса по сути является последовательным вызовом fork и execl (ещё можно между ними добавить dup2 или close), использован вместо fork и execl из-за похожести синтаксиса на CreateProcess, что повышает читаемость кода.
  \item kill - отправка процессу сигнала об окончании работы (SIGTERM или SIGKILL).
  \item pipe - создание pipe для POSIX.
  \item write - запись в поток.
  \item read - чтение из потока.
\end{itemize}
